\subsection{Puerta de Hadamard}
La puerta de Hadamard es una conocida puerta de un cúbit que transforma los estados de la base $\ket{0}$ y $\ket{1}$ en $\ket{+}=\frac{\ket{0}+\ket{1}}{\sqrt{2}}$ y $\ket{-}=\frac{\ket{0}-\ket{1}}{\sqrt{2}}$, respectivamente.

Así, la puerta de Hadamard puede expresarse usando la notación introducida como sigue:
\begin{equation}
  \label{eq:hadamard}
  \begin{split}
    \gate{H}_k\ket{a} & =a_{k0}\ket{a}^{k}_{+}+a_{k1}\ket{a}^{k}_{-}=                                                \\
                      & =\frac{a_{k0}+a_{k1}}{\sqrt{2}}\ket{a}^{k}_{0}+\frac{a_{k0}-a_{k1}}{\sqrt{2}}\ket{a}^{k}_{1}
  \end{split}
\end{equation}
De manera equivalente, como $a_{k+}=\braket{+}{a_k}=\frac{a_{k0}+a_{k1}}{\sqrt{2}}$
y $a_{k-}=\braket{-}{a_k}=\frac{a_{k0}-a_{k1}}{\sqrt{2}}$, la acción de
la puerta de Hadamard intercambia ambas bases:
\begin{equation*}
  \gate{H}_k\ket{a}=a_{k+}\ket{a}^{k}_{0}+a_{k-}\ket{a}^{k}_{1}
\end{equation*}

La conjugación por la puerta de Hadamard intercambia las puertas $\gate{X}$ y $\gate{Z}$ y cambia el signo de $\gate{Y}$:
\begin{align}
  \gate{H}_k\gate{X}_k\gate{H}_k & = \gate{Z}_k & \gate{H}_k\gate{Y}_k\gate{H}_k & = -\gate{Y}_k & \gate{H}_k\gate{Z}_k\gate{H}_k & = \gate{X}_k
\end{align}
